% --- Archon physics formalization source begin ---
% archon:physics
% archon:covers PhyXMiniProblems/problem_phyx_mini_0461.lean
% archon:source-report reports/phyx_mini/problem_phyx_mini_0461.source.json

\chapter{Physics problem phyx\_mini\_0461}
\label{ch:PhyXMiniProblems_problem_phyx_mini_0461}

\paragraph{Problem source.}
\#\# Physical scenario

An air pistol contains compressed air in a small cylinder, as shown in figure. Assume that the volume is \$1 \textbackslash{}, 	ext\{cm\}\textasciicircum{}3\$, the pressure is \$1 \textbackslash{}, \textbackslash{}text\{MPa\}\$, and the temperature is \$27 \textbackslash{}, \textasciicircum{}\{\textbackslash{}circ\}\textbackslash{}text\{C\}\$ when armed. A bullet, with \$m = 15 \textbackslash{}, \textbackslash{}text\{g\}\$, acts as a piston initially held by a pin (trigger); when released, the air expands in an isothermal process (\$T = 	ext\{constant\}\$). The air pressure is \$0.1 \textbackslash{}, \textbackslash{}text\{MPa\}\$ in the cylinder as the bullet leaves the gun.

\#\# Question

Find the the bullet exit velocity.

\#\# Answer choices

- A: A: 27.00 m/s

- B: B: 6.50 m/s

- C: C: 2.10 m/s

- D: D: 13.63 m/s

\#\# Figure caption (auxiliary; use the image as primary evidence)

The image depicts a diagram of a tube with various components:

1. **Tube and Sections**:
   - The tube is rectangular with a shaded outer boundary.
   - Inside the tube, there is a division between two sections: a blue section labeled "Air" and a hatched section.

2. **Air Section**:
   - The leftmost part of the tube is filled with a blue color and labeled "Air." This indicates this section contains air.

3. **Hatched Section**:
   - To the right of the air section, there is a hatched pattern that resembles the shape of a bullet or projectile, suggesting a separation or interface with the air section.

4. **Text and Notations**:
   - The text "Air" is present in the blue section.
   - To the right of the tube, there is a symbol \textbackslash{}( P\_0 \textbackslash{}), likely representing pressure outside the tube or in another section.

5. **Additional Elements**:
   - There is a curved element on the bottom left, resembling a lever or switch, which is connected to the internal division of the tube.

The diagram seems to illustrate a physical setup possibly related to gas dynamics or pressure systems within a closed environment.

\#\# Dataset metadata

Ideal Gases and Kinetic Theory; Physical Model Grounding Reasoning, Multi-Formula Reasoning

\paragraph{Recorded answer/context.}
D: D: 13.63 m/s

\paragraph{Figure/image path.}
/root/proposal\_for\_physic/science-mango/phyx\_mini\_run/phyx\_data/test\_image/461.png

\paragraph{Formalization target.}
create a compiling Lean file with sorry bodies at `PhyXMiniProblems/problem\_phyx\_mini\_0461.lean`. The Lean declarations must preserve the physical quantities, dimensions or dimensional roles, figure labels, governing-law hypotheses, and final relation expressed by this problem.
Use Mathlib/Physlib names found through LeanExplore where available. If a physics API is missing, introduce faithful local abstractions rather than scalar placeholder aliases.

\begin{theorem}[Physics formalization target]
\label{thm:physics:phyx_mini_0461:target}
\lean{PhyXMiniProblems.ProblemPhyXMini0461.problem_phyx_mini_0461}
\uses{def:physics:phyx-mini-0461:phyxminiproblems-problemphyxmini0461-speedinmeterspersecond, def:physics:phyx-mini-0461:phyxminiproblems-problemphyxmini0461-airpistolsetup, def:physics:phyx-mini-0461:phyxminiproblems-problemphyxmini0461-matchesproblemreadouts, def:physics:phyx-mini-0461:phyxminiproblems-problemphyxmini0461-matchesscenarioandsuppliedfigure, def:physics:phyx-mini-0461:phyxminiproblems-problemphyxmini0461-hasphysicalairpistolparameters, def:physics:phyx-mini-0461:phyxminiproblems-problemphyxmini0461-satisfiesisothermalidealgasexpansion, def:physics:phyx-mini-0461:phyxminiproblems-problemphyxmini0461-satisfiespressureworkrelations, def:physics:phyx-mini-0461:phyxminiproblems-problemphyxmini0461-satisfiesbulletworkenergylaw, def:physics:phyx-mini-0461:phyxminiproblems-problemphyxmini0461-roundstonearesthundredthmeterpersecond, def:physics:phyx-mini-0461:phyxminiproblems-problemphyxmini0461-matchesanswerchoice}
The assigned autoformalize agent should translate this physics problem into Lean declarations in the covered file, with theorem and lemma proof bodies written as `by sorry`.
\end{theorem}
\begin{proof}
This is an autoformalization task, not a proof task. Produce statements that can later be proved by the physics prover without weakening the physical model.
\end{proof}

% --- Archon declaration topology begin ---
\section{Declaration topology}
\begin{definition}[Volume Quantity]
  \label{def:physics:phyx-mini-0461:phyxminiproblems-problemphyxmini0461-volumequantity}
  \lean{PhyXMiniProblems.ProblemPhyXMini0461.VolumeQuantity}
  A nonnegative physical volume, with dimension length cubed
\end{definition}
\begin{proof}
  This is the declared model definition of Volume Quantity.
\end{proof}
\begin{definition}[Mass Quantity]
  \label{def:physics:phyx-mini-0461:phyxminiproblems-problemphyxmini0461-massquantity}
  \lean{PhyXMiniProblems.ProblemPhyXMini0461.MassQuantity}
  A nonnegative physical mass, independent of its readout unit
\end{definition}
\begin{proof}
  This is the declared model definition of Mass Quantity.
\end{proof}
\begin{definition}[Absolute Temperature Quantity]
  \label{def:physics:phyx-mini-0461:phyxminiproblems-problemphyxmini0461-absolutetemperaturequantity}
  \lean{PhyXMiniProblems.ProblemPhyXMini0461.AbsoluteTemperatureQuantity}
  A nonnegative absolute temperature, independent of its readout unit
\end{definition}
\begin{proof}
  This is the declared model definition of Absolute Temperature Quantity.
\end{proof}
\begin{definition}[Pressure Quantity]
  \label{def:physics:phyx-mini-0461:phyxminiproblems-problemphyxmini0461-pressurequantity}
  \lean{PhyXMiniProblems.ProblemPhyXMini0461.PressureQuantity}
  A physical pressure
\end{definition}
\begin{proof}
  This is the declared model definition of Pressure Quantity.
\end{proof}
\begin{definition}[Speed Quantity]
  \label{def:physics:phyx-mini-0461:phyxminiproblems-problemphyxmini0461-speedquantity}
  \lean{PhyXMiniProblems.ProblemPhyXMini0461.SpeedQuantity}
  A nonnegative physical speed
\end{definition}
\begin{proof}
  This is the declared model definition of Speed Quantity.
\end{proof}
\begin{definition}[Energy Quantity]
  \label{def:physics:phyx-mini-0461:phyxminiproblems-problemphyxmini0461-energyquantity}
  \lean{PhyXMiniProblems.ProblemPhyXMini0461.EnergyQuantity}
  A physical work or energy
\end{definition}
\begin{proof}
  This is the declared model definition of Energy Quantity.
\end{proof}
\begin{definition}[volume Readout]
  \label{def:physics:phyx-mini-0461:phyxminiproblems-problemphyxmini0461-volumereadout}
  \lean{PhyXMiniProblems.ProblemPhyXMini0461.volumeReadout}
  \uses{def:physics:phyx-mini-0461:phyxminiproblems-problemphyxmini0461-volumequantity}
  Read a volume in the cube of the selected length unit
\end{definition}
\begin{proof}
  This is the declared model definition of volume Readout.
\end{proof}
\begin{definition}[mass Readout]
  \label{def:physics:phyx-mini-0461:phyxminiproblems-problemphyxmini0461-massreadout}
  \lean{PhyXMiniProblems.ProblemPhyXMini0461.massReadout}
  \uses{def:physics:phyx-mini-0461:phyxminiproblems-problemphyxmini0461-massquantity}
  Read a mass in the selected mass unit
\end{definition}
\begin{proof}
  This is the declared model definition of mass Readout.
\end{proof}
\begin{definition}[absolute Temperature Readout]
  \label{def:physics:phyx-mini-0461:phyxminiproblems-problemphyxmini0461-absolutetemperaturereadout}
  \lean{PhyXMiniProblems.ProblemPhyXMini0461.absoluteTemperatureReadout}
  \uses{def:physics:phyx-mini-0461:phyxminiproblems-problemphyxmini0461-absolutetemperaturequantity}
  Read an absolute temperature in the selected temperature unit
\end{definition}
\begin{proof}
  This is the declared model definition of absolute Temperature Readout.
\end{proof}
\begin{definition}[pressure In Pascals]
  \label{def:physics:phyx-mini-0461:phyxminiproblems-problemphyxmini0461-pressureinpascals}
  \lean{PhyXMiniProblems.ProblemPhyXMini0461.pressureInPascals}
  \uses{def:physics:phyx-mini-0461:phyxminiproblems-problemphyxmini0461-pressurequantity}
  Pressure readout in pascals
\end{definition}
\begin{proof}
  This is the declared model definition of pressure In Pascals.
\end{proof}
\begin{definition}[pressure In Megapascals]
  \label{def:physics:phyx-mini-0461:phyxminiproblems-problemphyxmini0461-pressureinmegapascals}
  \lean{PhyXMiniProblems.ProblemPhyXMini0461.pressureInMegapascals}
  \uses{def:physics:phyx-mini-0461:phyxminiproblems-problemphyxmini0461-pressurequantity, def:physics:phyx-mini-0461:phyxminiproblems-problemphyxmini0461-pressureinpascals}
  Pressure readout in megapascals
\end{definition}
\begin{proof}
  This is the declared model definition of pressure In Megapascals.
\end{proof}
\begin{definition}[volume In Cubic Meters]
  \label{def:physics:phyx-mini-0461:phyxminiproblems-problemphyxmini0461-volumeincubicmeters}
  \lean{PhyXMiniProblems.ProblemPhyXMini0461.volumeInCubicMeters}
  \uses{def:physics:phyx-mini-0461:phyxminiproblems-problemphyxmini0461-volumequantity, def:physics:phyx-mini-0461:phyxminiproblems-problemphyxmini0461-volumereadout}
  Volume readout in cubic meters
\end{definition}
\begin{proof}
  This is the declared model definition of volume In Cubic Meters.
\end{proof}
\begin{definition}[volume In Cubic Centimeters]
  \label{def:physics:phyx-mini-0461:phyxminiproblems-problemphyxmini0461-volumeincubiccentimeters}
  \lean{PhyXMiniProblems.ProblemPhyXMini0461.volumeInCubicCentimeters}
  \uses{def:physics:phyx-mini-0461:phyxminiproblems-problemphyxmini0461-volumequantity, def:physics:phyx-mini-0461:phyxminiproblems-problemphyxmini0461-volumereadout}
  Volume readout in cubic centimeters
\end{definition}
\begin{proof}
  This is the declared model definition of volume In Cubic Centimeters.
\end{proof}
\begin{definition}[mass In Kilograms]
  \label{def:physics:phyx-mini-0461:phyxminiproblems-problemphyxmini0461-massinkilograms}
  \lean{PhyXMiniProblems.ProblemPhyXMini0461.massInKilograms}
  \uses{def:physics:phyx-mini-0461:phyxminiproblems-problemphyxmini0461-massquantity, def:physics:phyx-mini-0461:phyxminiproblems-problemphyxmini0461-massreadout}
  Mass readout in kilograms
\end{definition}
\begin{proof}
  This is the declared model definition of mass In Kilograms.
\end{proof}
\begin{definition}[mass In Grams]
  \label{def:physics:phyx-mini-0461:phyxminiproblems-problemphyxmini0461-massingrams}
  \lean{PhyXMiniProblems.ProblemPhyXMini0461.massInGrams}
  \uses{def:physics:phyx-mini-0461:phyxminiproblems-problemphyxmini0461-massquantity, def:physics:phyx-mini-0461:phyxminiproblems-problemphyxmini0461-massreadout}
  Mass readout in grams
\end{definition}
\begin{proof}
  This is the declared model definition of mass In Grams.
\end{proof}
\begin{definition}[temperature In Kelvin]
  \label{def:physics:phyx-mini-0461:phyxminiproblems-problemphyxmini0461-temperatureinkelvin}
  \lean{PhyXMiniProblems.ProblemPhyXMini0461.temperatureInKelvin}
  \uses{def:physics:phyx-mini-0461:phyxminiproblems-problemphyxmini0461-absolutetemperaturequantity, def:physics:phyx-mini-0461:phyxminiproblems-problemphyxmini0461-absolutetemperaturereadout}
  Absolute-temperature readout in kelvin
\end{definition}
\begin{proof}
  This is the declared model definition of temperature In Kelvin.
\end{proof}
\begin{definition}[temperature In Degrees Celsius]
  \label{def:physics:phyx-mini-0461:phyxminiproblems-problemphyxmini0461-temperatureindegreescelsius}
  \lean{PhyXMiniProblems.ProblemPhyXMini0461.temperatureInDegreesCelsius}
  \uses{def:physics:phyx-mini-0461:phyxminiproblems-problemphyxmini0461-absolutetemperaturequantity, def:physics:phyx-mini-0461:phyxminiproblems-problemphyxmini0461-temperatureinkelvin}
  Celsius readout obtained from the absolute Kelvin readout
\end{definition}
\begin{proof}
  This is the declared model definition of temperature In Degrees Celsius.
\end{proof}
\begin{definition}[speed In Meters Per Second]
  \label{def:physics:phyx-mini-0461:phyxminiproblems-problemphyxmini0461-speedinmeterspersecond}
  \lean{PhyXMiniProblems.ProblemPhyXMini0461.speedInMetersPerSecond}
  \uses{def:physics:phyx-mini-0461:phyxminiproblems-problemphyxmini0461-speedquantity}
  Speed readout in meters per second
\end{definition}
\begin{proof}
  This is the declared model definition of speed In Meters Per Second.
\end{proof}
\begin{definition}[energy In Joules]
  \label{def:physics:phyx-mini-0461:phyxminiproblems-problemphyxmini0461-energyinjoules}
  \lean{PhyXMiniProblems.ProblemPhyXMini0461.energyInJoules}
  \uses{def:physics:phyx-mini-0461:phyxminiproblems-problemphyxmini0461-energyquantity}
  Energy readout in joules
\end{definition}
\begin{proof}
  This is the declared model definition of energy In Joules.
\end{proof}
\begin{definition}[Air State]
  \label{def:physics:phyx-mini-0461:phyxminiproblems-problemphyxmini0461-airstate}
  \lean{PhyXMiniProblems.ProblemPhyXMini0461.AirState}
  \uses{def:physics:phyx-mini-0461:phyxminiproblems-problemphyxmini0461-volumequantity, def:physics:phyx-mini-0461:phyxminiproblems-problemphyxmini0461-absolutetemperaturequantity, def:physics:phyx-mini-0461:phyxminiproblems-problemphyxmini0461-pressurequantity}
  Thermodynamic data for the air at one stage of the shot
\end{definition}
\begin{proof}
  This is the declared model definition of Air State.
\end{proof}
\begin{definition}[Figure Component]
  \label{def:physics:phyx-mini-0461:phyxminiproblems-problemphyxmini0461-figurecomponent}
  \lean{PhyXMiniProblems.ProblemPhyXMini0461.FigureComponent}
  Components distinguishable in the supplied air-pistol figure
\end{definition}
\begin{proof}
  This is the declared model definition of Figure Component.
\end{proof}
\begin{definition}[Figure Label]
  \label{def:physics:phyx-mini-0461:phyxminiproblems-problemphyxmini0461-figurelabel}
  \lean{PhyXMiniProblems.ProblemPhyXMini0461.FigureLabel}
  Text or mathematical labels printed in the supplied figure
\end{definition}
\begin{proof}
  This is the declared model definition of Figure Label.
\end{proof}
\begin{definition}[Thermodynamic Process]
  \label{def:physics:phyx-mini-0461:phyxminiproblems-problemphyxmini0461-thermodynamicprocess}
  \lean{PhyXMiniProblems.ProblemPhyXMini0461.ThermodynamicProcess}
  The thermodynamic-process idealization stated in the problem
\end{definition}
\begin{proof}
  This is the declared model definition of Thermodynamic Process.
\end{proof}
\begin{definition}[Bullet Role]
  \label{def:physics:phyx-mini-0461:phyxminiproblems-problemphyxmini0461-bulletrole}
  \lean{PhyXMiniProblems.ProblemPhyXMini0461.BulletRole}
  Mechanical role played by the bullet in the air chamber
\end{definition}
\begin{proof}
  This is the declared model definition of Bullet Role.
\end{proof}
\begin{definition}[Release Mechanism]
  \label{def:physics:phyx-mini-0461:phyxminiproblems-problemphyxmini0461-releasemechanism}
  \lean{PhyXMiniProblems.ProblemPhyXMini0461.ReleaseMechanism}
  Release mechanism shown at the lower side of the barrel
\end{definition}
\begin{proof}
  This is the declared model definition of Release Mechanism.
\end{proof}
\begin{definition}[Supplied Air Pistol Figure]
  \label{def:physics:phyx-mini-0461:phyxminiproblems-problemphyxmini0461-suppliedairpistolfigure}
  \lean{PhyXMiniProblems.ProblemPhyXMini0461.SuppliedAirPistolFigure}
  \uses{def:physics:phyx-mini-0461:phyxminiproblems-problemphyxmini0461-pressurequantity, def:physics:phyx-mini-0461:phyxminiproblems-problemphyxmini0461-figurecomponent, def:physics:phyx-mini-0461:phyxminiproblems-problemphyxmini0461-figurelabel}
  Qualitative evidence transcribed from the primary image. In particular, the hatched bullet separates the blue air region from the open bore, the pin initially retains it, and the exterior pressure is labelled P₀
\end{definition}
\begin{proof}
  This is the declared model definition of Supplied Air Pistol Figure.
\end{proof}
\begin{definition}[Air Pistol Setup]
  \label{def:physics:phyx-mini-0461:phyxminiproblems-problemphyxmini0461-airpistolsetup}
  \lean{PhyXMiniProblems.ProblemPhyXMini0461.AirPistolSetup}
  \uses{def:physics:phyx-mini-0461:phyxminiproblems-problemphyxmini0461-massquantity, def:physics:phyx-mini-0461:phyxminiproblems-problemphyxmini0461-pressurequantity, def:physics:phyx-mini-0461:phyxminiproblems-problemphyxmini0461-speedquantity, def:physics:phyx-mini-0461:phyxminiproblems-problemphyxmini0461-energyquantity, def:physics:phyx-mini-0461:phyxminiproblems-problemphyxmini0461-airstate, def:physics:phyx-mini-0461:phyxminiproblems-problemphyxmini0461-thermodynamicprocess, def:physics:phyx-mini-0461:phyxminiproblems-problemphyxmini0461-bulletrole, def:physics:phyx-mini-0461:phyxminiproblems-problemphyxmini0461-releasemechanism, def:physics:phyx-mini-0461:phyxminiproblems-problemphyxmini0461-suppliedairpistolfigure}
  All physical quantities needed by the model. The exit speed is an unconstrained physical speed here; no structure field assigns it an answer choice or numerical value
\end{definition}
\begin{proof}
  This is the declared model definition of Air Pistol Setup.
\end{proof}
\begin{definition}[Matches Problem Readouts]
  \label{def:physics:phyx-mini-0461:phyxminiproblems-problemphyxmini0461-matchesproblemreadouts}
  \lean{PhyXMiniProblems.ProblemPhyXMini0461.MatchesProblemReadouts}
  \uses{def:physics:phyx-mini-0461:phyxminiproblems-problemphyxmini0461-pressureinmegapascals, def:physics:phyx-mini-0461:phyxminiproblems-problemphyxmini0461-volumeincubiccentimeters, def:physics:phyx-mini-0461:phyxminiproblems-problemphyxmini0461-massingrams, def:physics:phyx-mini-0461:phyxminiproblems-problemphyxmini0461-temperatureindegreescelsius, def:physics:phyx-mini-0461:phyxminiproblems-problemphyxmini0461-speedinmeterspersecond, def:physics:phyx-mini-0461:phyxminiproblems-problemphyxmini0461-airpistolsetup}
  Numerical readouts stated in the problem. The 0.101 MPa value for P₀ is the usual textbook standard-atmosphere approximation implicit in the figure and is the value used by the recorded answer. No exit volume, work, or exit-speed value occurs among these premises
\end{definition}
\begin{proof}
  This is the declared model definition of Matches Problem Readouts.
\end{proof}
\begin{definition}[Matches Scenario And Supplied Figure]
  \label{def:physics:phyx-mini-0461:phyxminiproblems-problemphyxmini0461-matchesscenarioandsuppliedfigure}
  \lean{PhyXMiniProblems.ProblemPhyXMini0461.MatchesScenarioAndSuppliedFigure}
  \uses{def:physics:phyx-mini-0461:phyxminiproblems-problemphyxmini0461-airpistolsetup}
  Qualitative prose assumptions together with the primary-image facts
\end{definition}
\begin{proof}
  This is the declared model definition of Matches Scenario And Supplied Figure.
\end{proof}
\begin{definition}[Has Physical Air Pistol Parameters]
  \label{def:physics:phyx-mini-0461:phyxminiproblems-problemphyxmini0461-hasphysicalairpistolparameters}
  \lean{PhyXMiniProblems.ProblemPhyXMini0461.HasPhysicalAirPistolParameters}
  \uses{def:physics:phyx-mini-0461:phyxminiproblems-problemphyxmini0461-pressureinpascals, def:physics:phyx-mini-0461:phyxminiproblems-problemphyxmini0461-volumeincubicmeters, def:physics:phyx-mini-0461:phyxminiproblems-problemphyxmini0461-massinkilograms, def:physics:phyx-mini-0461:phyxminiproblems-problemphyxmini0461-temperatureinkelvin, def:physics:phyx-mini-0461:phyxminiproblems-problemphyxmini0461-speedinmeterspersecond, def:physics:phyx-mini-0461:phyxminiproblems-problemphyxmini0461-airpistolsetup}
  Positivity and nondegeneracy conditions for the physical shot
\end{definition}
\begin{proof}
  This is the declared model definition of Has Physical Air Pistol Parameters.
\end{proof}
\begin{definition}[Satisfies Isothermal Ideal Gas Expansion]
  \label{def:physics:phyx-mini-0461:phyxminiproblems-problemphyxmini0461-satisfiesisothermalidealgasexpansion}
  \lean{PhyXMiniProblems.ProblemPhyXMini0461.SatisfiesIsothermalIdealGasExpansion}
  \uses{def:physics:phyx-mini-0461:phyxminiproblems-problemphyxmini0461-airpistolsetup}
  The fixed amount of ideal gas obeys P V = constant during the expansion, and its absolute temperature is unchanged. Both statements are required in every coherent unit system. Neither law assigns a value to the bullet speed
\end{definition}
\begin{proof}
  This is the declared model definition of Satisfies Isothermal Ideal Gas Expansion.
\end{proof}
\begin{definition}[Satisfies Pressure Work Relations]
  \label{def:physics:phyx-mini-0461:phyxminiproblems-problemphyxmini0461-satisfiespressureworkrelations}
  \lean{PhyXMiniProblems.ProblemPhyXMini0461.SatisfiesPressureWorkRelations}
  \uses{def:physics:phyx-mini-0461:phyxminiproblems-problemphyxmini0461-airpistolsetup}
  For an isothermal ideal-gas expansion, the gas boundary work is Pᵢ Vᵢ log (V\_f / Vᵢ). Motion against the constant exterior pressure does work P₀ (V\_f - Vᵢ). The equations are stated in every coherent unit system, so each side has the readout dimension of energy
\end{definition}
\begin{proof}
  This is the declared model definition of Satisfies Pressure Work Relations.
\end{proof}
\begin{definition}[Satisfies Bullet Work Energy Law]
  \label{def:physics:phyx-mini-0461:phyxminiproblems-problemphyxmini0461-satisfiesbulletworkenergylaw}
  \lean{PhyXMiniProblems.ProblemPhyXMini0461.SatisfiesBulletWorkEnergyLaw}
  \uses{def:physics:phyx-mini-0461:phyxminiproblems-problemphyxmini0461-airpistolsetup}
  With friction and other losses neglected, the net pressure work changes the bullet's translational kinetic energy. The initial kinetic term is retained explicitly rather than hidden by the rest readout
\end{definition}
\begin{proof}
  This is the declared model definition of Satisfies Bullet Work Energy Law.
\end{proof}
\begin{lemma}[exit Volume In Cubic Centimeters eq ten]
  \label{lem:physics:phyx-mini-0461:phyxminiproblems-problemphyxmini0461-exitvolumeincubiccentimeters-eq-ten}
  \lean{PhyXMiniProblems.ProblemPhyXMini0461.exitVolumeInCubicCentimeters_eq_ten}
  \uses{def:physics:phyx-mini-0461:phyxminiproblems-problemphyxmini0461-volumeincubiccentimeters, def:physics:phyx-mini-0461:phyxminiproblems-problemphyxmini0461-airpistolsetup, def:physics:phyx-mini-0461:phyxminiproblems-problemphyxmini0461-matchesproblemreadouts, def:physics:phyx-mini-0461:phyxminiproblems-problemphyxmini0461-hasphysicalairpistolparameters, def:physics:phyx-mini-0461:phyxminiproblems-problemphyxmini0461-satisfiesisothermalidealgasexpansion}
  The ideal-gas invariant determines an exit volume of 10 cm³
\end{lemma}
\begin{proof}
  The conclusion follows from the governing relations and figure readouts named in the statement of exit Volume In Cubic Centimeters eq ten.
\end{proof}
\begin{lemma}[gas Boundary Work In Joules eq log ten]
  \label{lem:physics:phyx-mini-0461:phyxminiproblems-problemphyxmini0461-gasboundaryworkinjoules-eq-log-ten}
  \lean{PhyXMiniProblems.ProblemPhyXMini0461.gasBoundaryWorkInJoules_eq_log_ten}
  \uses{def:physics:phyx-mini-0461:phyxminiproblems-problemphyxmini0461-energyinjoules, def:physics:phyx-mini-0461:phyxminiproblems-problemphyxmini0461-airpistolsetup, def:physics:phyx-mini-0461:phyxminiproblems-problemphyxmini0461-matchesproblemreadouts, def:physics:phyx-mini-0461:phyxminiproblems-problemphyxmini0461-hasphysicalairpistolparameters, def:physics:phyx-mini-0461:phyxminiproblems-problemphyxmini0461-satisfiesisothermalidealgasexpansion, def:physics:phyx-mini-0461:phyxminiproblems-problemphyxmini0461-satisfiespressureworkrelations}
  The isothermal gas does log 10 joules of positive boundary work
\end{lemma}
\begin{proof}
  The conclusion follows from the governing relations and figure readouts named in the statement of gas Boundary Work In Joules eq log ten.
\end{proof}
\begin{lemma}[ambient Displacement Work In Joules eq]
  \label{lem:physics:phyx-mini-0461:phyxminiproblems-problemphyxmini0461-ambientdisplacementworkinjoules-eq}
  \lean{PhyXMiniProblems.ProblemPhyXMini0461.ambientDisplacementWorkInJoules_eq}
  \uses{def:physics:phyx-mini-0461:phyxminiproblems-problemphyxmini0461-energyinjoules, def:physics:phyx-mini-0461:phyxminiproblems-problemphyxmini0461-airpistolsetup, def:physics:phyx-mini-0461:phyxminiproblems-problemphyxmini0461-matchesproblemreadouts, def:physics:phyx-mini-0461:phyxminiproblems-problemphyxmini0461-hasphysicalairpistolparameters, def:physics:phyx-mini-0461:phyxminiproblems-problemphyxmini0461-satisfiesisothermalidealgasexpansion, def:physics:phyx-mini-0461:phyxminiproblems-problemphyxmini0461-satisfiespressureworkrelations}
  At P₀ = 0.101 MPa, displacing 9 cm³ costs 0.909 J
\end{lemma}
\begin{proof}
  The conclusion follows from the governing relations and figure readouts named in the statement of ambient Displacement Work In Joules eq.
\end{proof}
\begin{lemma}[bullet Exit Speed exact readout]
  \label{lem:physics:phyx-mini-0461:phyxminiproblems-problemphyxmini0461-bulletexitspeed-exact-readout}
  \lean{PhyXMiniProblems.ProblemPhyXMini0461.bulletExitSpeed_exact_readout}
  \uses{def:physics:phyx-mini-0461:phyxminiproblems-problemphyxmini0461-speedinmeterspersecond, def:physics:phyx-mini-0461:phyxminiproblems-problemphyxmini0461-airpistolsetup, def:physics:phyx-mini-0461:phyxminiproblems-problemphyxmini0461-matchesproblemreadouts, def:physics:phyx-mini-0461:phyxminiproblems-problemphyxmini0461-hasphysicalairpistolparameters, def:physics:phyx-mini-0461:phyxminiproblems-problemphyxmini0461-satisfiesisothermalidealgasexpansion, def:physics:phyx-mini-0461:phyxminiproblems-problemphyxmini0461-satisfiespressureworkrelations, def:physics:phyx-mini-0461:phyxminiproblems-problemphyxmini0461-satisfiesbulletworkenergylaw}
  The positive solution of the work-energy equation. This exact radical is a derived conclusion; the decimal answer is obtained only by rounding it
\end{lemma}
\begin{proof}
  The conclusion follows from the governing relations and figure readouts named in the statement of bullet Exit Speed exact readout.
\end{proof}
\begin{definition}[Answer Choice]
  \label{def:physics:phyx-mini-0461:phyxminiproblems-problemphyxmini0461-answerchoice}
  \lean{PhyXMiniProblems.ProblemPhyXMini0461.AnswerChoice}
  Labels of the four displayed answer choices
\end{definition}
\begin{proof}
  This is the declared model definition of Answer Choice.
\end{proof}
\begin{definition}[displayed Answer Speed In Meters Per Second]
  \label{def:physics:phyx-mini-0461:phyxminiproblems-problemphyxmini0461-displayedanswerspeedinmeterspersecond}
  \lean{PhyXMiniProblems.ProblemPhyXMini0461.displayedAnswerSpeedInMetersPerSecond}
  \uses{def:physics:phyx-mini-0461:phyxminiproblems-problemphyxmini0461-answerchoice}
  Speed in meters per second printed beside each answer choice
\end{definition}
\begin{proof}
  This is the declared model definition of displayed Answer Speed In Meters Per Second.
\end{proof}
\begin{definition}[recorded Dataset Answer]
  \label{def:physics:phyx-mini-0461:phyxminiproblems-problemphyxmini0461-recordeddatasetanswer}
  \lean{PhyXMiniProblems.ProblemPhyXMini0461.recordedDatasetAnswer}
  \uses{def:physics:phyx-mini-0461:phyxminiproblems-problemphyxmini0461-answerchoice}
  The answer label recorded in the supplied dataset; it is not a premise
\end{definition}
\begin{proof}
  This is the declared model definition of recorded Dataset Answer.
\end{proof}
\begin{definition}[Rounds To Nearest Hundredth Meter Per Second]
  \label{def:physics:phyx-mini-0461:phyxminiproblems-problemphyxmini0461-roundstonearesthundredthmeterpersecond}
  \lean{PhyXMiniProblems.ProblemPhyXMini0461.RoundsToNearestHundredthMeterPerSecond}
  \uses{def:physics:phyx-mini-0461:phyxminiproblems-problemphyxmini0461-speedquantity, def:physics:phyx-mini-0461:phyxminiproblems-problemphyxmini0461-speedinmeterspersecond}
  Agreement with a speed printed to the nearest hundredth meter per second
\end{definition}
\begin{proof}
  This is the declared model definition of Rounds To Nearest Hundredth Meter Per Second.
\end{proof}
\begin{definition}[Matches Answer Choice]
  \label{def:physics:phyx-mini-0461:phyxminiproblems-problemphyxmini0461-matchesanswerchoice}
  \lean{PhyXMiniProblems.ProblemPhyXMini0461.MatchesAnswerChoice}
  \uses{def:physics:phyx-mini-0461:phyxminiproblems-problemphyxmini0461-airpistolsetup, def:physics:phyx-mini-0461:phyxminiproblems-problemphyxmini0461-answerchoice, def:physics:phyx-mini-0461:phyxminiproblems-problemphyxmini0461-displayedanswerspeedinmeterspersecond, def:physics:phyx-mini-0461:phyxminiproblems-problemphyxmini0461-roundstonearesthundredthmeterpersecond}
  The physical exit speed rounds to the number printed beside a choice
\end{definition}
\begin{proof}
  This is the declared model definition of Matches Answer Choice.
\end{proof}
% --- Archon declaration topology end ---
% --- Archon physics formalization source end ---
